% ===== 原书 p.46（下半页，附录 B 起） =====
\section{Atiyah--Hirzebruch 谱序列}

这里给出计算 K 理论的 Atiyah--Hirzebruch 谱序列，叙述依 [8, §2]。我们不假定读者熟悉由 Cartan--Eilenberg 系统构造谱序列（见 [14]），而是从对与三元组的长正合列出发，推出所需的全部同调代数。

设 $X$ 是有限（连通）CW 复形，$X^p$ 记它的 $p$-骨架。定义
\[
K^*_p(X)=\Ker\{K^*(X)\to K^*(X^{p-1})\}.
\]
这就是 K 理论的 Atiyah--Hirzebruch 滤过，且
\[
K^*_0(X)=K^*(X)\supset K^*_1(X)=\tilde K^*(X)\supset\cdots\supset K^*_{\max}(X)=0,
\]
其中 $X^{-1}=\varnothing$，$X^{\max-1}=X$。\footnote{尽管从定义看不出，Atiyah--Hirzebruch 滤过其实是同伦型不变量，与 $X$ 上 CW 结构的选取无关。此外这一滤过是乘性的，即 $K^*_p(X)\cdot K^*_q(X)\subset K^*_{p+q}(X)$ 对一切 $p,q$ 成立。更多信息见 [8, §2]。} Atiyah--Hirzebruch 谱序列计算相伴分次分量
\[
\operatorname{Gr}_p K^*(X)=K^*_p(X)\big/K^*_{p+1}(X).
\]
若知道这些分次分量，就可以用一串群扩张去还原 $K^*(X)$。

现在考虑如下正合列的交换图表：
\[
\begin{tikzcd}[column sep=small, row sep=large]
K^{*-1}(X^p) \arrow[r] \arrow[d] & K^*(X,X^p) \arrow[r] \arrow[d] & K^*(X) \arrow[r] \arrow[d,equal] & K^*(X^p) \arrow[d] \\
K^{*-1}(X^{p-1}) \arrow[r] \arrow[d] & K^*(X,X^{p-1}) \arrow[r] \arrow[d] & K^*(X) \arrow[r] & K^*(X^{p-1}) \\
K^*(X^p,X^{p-1}) \arrow[r,equal] \arrow[d] & K^*(X^p,X^{p-1}) \arrow[d] & & \\
K^*(X^p) \arrow[r] & K^{*+1}(X,X^p) & &
\end{tikzcd}
\]
其中上两行是 $(X,X^p)$ 与 $(X,X^{p-1})$ 的长正合列，左两列是 $(X^p,X^{p-1})$ 与三元组 $(X,X^p,X^{p-1})$ 的长正合列。

\transnote{图表中两处双线含义不同：$K^*(X)$ 处的双线只是把上两行的同一列对齐，标出两行长正合列在此共用该群；$K^*(X^p,X^{p-1})$ 处的双线则表示同一个相对群。两者都不是同构箭头，读图时不必找对应的映射。}

% ===== 原书 p.47 =====
由 $K^*_p(X)$ 的定义与上图表上两行的正合性，在 $K^*(X)$ 中有
\[
\begin{aligned}
K^*_p(X)&=\operatorname{Im} K^*(X,X^{p-1})\iso K^*(X,X^{p-1})\big/\operatorname{Im} K^{*-1}(X^{p-1}),\\
K^*_{p+1}(X)&=\operatorname{Im} K^*(X,X^p)\iso K^*(X,X^p)\big/\operatorname{Im} K^{*-1}(X^p),
\end{aligned}
\]
于是相伴分次分量为
\[
\begin{aligned}
\operatorname{Gr}_p K^*(X)&\iso K^*(X,X^{p-1})\big/\bigl(\operatorname{Im} K^*(X,X^p)+\operatorname{Im} K^{*-1}(X^{p-1})\bigr)\\
&\iso\bigl(K^*(X,X^{p-1})\big/\operatorname{Im} K^*(X,X^p)\bigr)\big/\operatorname{Im} K^{*-1}(X^{p-1})\\
&\iso\Ker\{K^*(X^p,X^{p-1})\to K^{*+1}(X,X^p)\}\big/\operatorname{Im} K^{*-1}(X^{p-1}).
\end{aligned}
\]
为计算这个分次分量，构造一个谱序列，其第一页为
\[
E^p_1:=K^*(X^p,X^{p-1}),
\]
第 $k$ 页由
\[
E^p_k=\frac{\Ker\{K^*(X^p,X^{p-1})\xrightarrow{\ \beta\ }K^{*+1}(X^{p-1+k},X^p)\}}
{\operatorname{Im}\{K^{*-1}(X^{p-1},X^{p-k})\xrightarrow{\ \beta'\ }K^*(X^p,X^{p-1})\}}
\]
给出。注意对 $E_1$，从
\[
E^p_1=\Ker\{K^*(X^p,X^{p-1})\to K^{*+1}(X^p,X^p)\}\big/\operatorname{Im} K^{*-1}(X^{p-1},X^{p-1})
\]
出发，这与 $E^p_1$ 的定义一致。另一方面，$k$ 足够大时 $X^{p-1+k}=X$、$X^{p-k}=\varnothing$，于是序列稳定下来，给出所需的分次分量 $E^p_\infty=\operatorname{Gr}_p K^*(X)$。考察下面的交换图表（取自前一图表的左下角），
\[
\begin{tikzcd}[row sep=large, column sep=large]
K^{*-1}(X^{p-1},X^{p-k}) \arrow[r] \arrow[d] & K^*(X^{p-1+k},X^{p-1}) \arrow[d] \\
K^*(X^p,X^{p-1}) \arrow[r,equal] & K^*(X^p,X^{p-1}) \arrow[d] \\
& K^{*+1}(X^{p-1+k},X^p)
\end{tikzcd}
\]
可见在 $K^*(X^p,X^{p-1})$ 中
\[
\begin{aligned}
\operatorname{Im} K^{*-1}(X^{p-1},X^{p-k})&\subset\operatorname{Im} K^{*-1}(X^{p-1+k},X^{p-1})\\
&=\Ker\{K^*(X^p,X^{p-1})\to K^{*+1}(X^{p-1+k},X^p)\},
\end{aligned}
\]
从而 $E^p_k$ 确实良定义。

\transnote{稳定要两个条件同时成立：$X^{p-1+k}=X$ 需 $p-1+k\ge\max$，$X^{p-k}=\varnothing$ 需 $p-k\le-1$。取 $k\ge\max-p+1$ 即可，故每个 $p$ 上 $E^p_k$ 终将稳定，$\max$ 是最高非空骨架的序号。}

% ===== 原书 p.48 =====
为构造谱序列的微分，考虑如下交换图表：
\[
\begin{tikzcd}[row sep=large, column sep=small]
& K^{*-1}(X^{p-1},X^{p-k}) \arrow[r] \arrow[dr,"\beta'"'] & K^{*-1}(X^{p-1},X^{p-k-1}) \arrow[d,"\alpha'"] & \\
& & K^*(X^p,X^{p-1}) \arrow[dll,dashed,"d_k"'] \arrow[d,"\alpha"] \arrow[dr,"\beta"] & \\
K^*(X^{p-1+k},X^p) \arrow[r,"\epsilon"] & K^{*+1}(X^{p+k},X^{p-1+k}) \arrow[r,"\delta"] & K^{*+1}(X^{p+k},X^p) \arrow[r,"\gamma"] & K^{*+1}(X^{p-1+k},X^p)
\end{tikzcd}
\]
其中上行来自三元组 $(X^{p-1},X^{p-k},X^{p-k-1})$ 的长正合列，下行来自三元组 $(X^{p+k},X^{p-1+k},X^p)$ 的长正合列。目标是：用群 $E^p_k=\Ker\beta/\operatorname{Im}\beta'$ 计算群 $E^{p+k}_{k+1}=\Ker\alpha/\operatorname{Im}\alpha'$。

设 $x\in K^*(X^p,X^{p-1})$ 是类 $[x]\in E^p_k$ 的代表元。由 $E^p_k$ 的定义 $x\in\Ker\beta$，于是由交换图表 $\alpha x\in\Ker\gamma=\operatorname{Im}\delta$。因此可写 $\alpha x=\delta y$，其中 $y\in K^{*+1}(X^{p+k},X^{p-1+k})$，且 $y$ 在 $\Ker\delta=\operatorname{Im}\epsilon$ 的意义下确定。这样就得到良定义的类
\[
[y]\in K^{*+1}(X^{p+k},X^{p-1+k})\Big/\operatorname{Im}\Bigl\{K^*(X^{p-1+k},X^p)\xrightarrow{\ \epsilon\ }K^{*+1}(X^{p+K},X^{p-1+k})\Bigr\},
\]
并把谱序列的微分定义为
\[
d_k:E^p_k\to E^{p+k}_k,\qquad d_k[x]:=[y].
\]
为完备起见，还应验证
\[
y\in\Ker\{K^{*+1}(X^{p+k},X^{p-1+k})\to K^{*+2}(X^{p+2k-1},X^{p+k})\}
\]
（用类似的图表追踪），验证类 $d_k[x]\in E^{p+k,*+1}_k$ 与 $x$ 的选取无关（因为 $\operatorname{Im}\beta'\subset\operatorname{Im}\alpha'\subset\Ker\alpha$），并验证 $d_k$-上同调确实给出
\[
\Ker d_k/\operatorname{Im} d_k=\Ker\alpha/\operatorname{Im}\alpha'=E^p_{k+1},
\]
即谱序列的下一页。

为把上同调分次并入上述讨论，在谱序列上使用双分次：项 $E^{p,q}_k$ 带的上同调次数为 $p+q$。微分于是给出
\[
d_k:E^{p,q}_k\to E^{p+k,q-k+1}_k
\]
对一切 $p,q$ 成立。这个谱序列的 $E_1$ 页有分量 $E^{p,q}_1=K^{p+q}(X^p,X^{p-1})$。然而，若相对 $(p-1)$-骨架来作，则 $X^p$ 成为一束 $p$-球面（bouquet），每个 $p$-胞腔对应一个。于是得到
\[
E^{p,q}_1=K^{p+q}(X^p,X^{p-1})\iso K^{p+q}(S^p)^{\oplus\beta_p}\iso K^q(\pt)^{\oplus\beta_p}=C^p\bigl(X;K^q(\pt)\bigr)
\]
（其中 $\beta_p$ 是 $p$-胞腔的个数），这就是 $X$ 的、以一点的 K 理论为系数的胞腔上同调的链群。此外，第一个微分
\[
d_1:C^p\bigl(X;K^q(\pt)\bigr)\to C^{p+1}\bigl(X;K^q(\pt)\bigr)
\]

% ===== 原书 p.49 =====
正是胞腔链复形中的微分（见 [2, §III.7]），于是
\[
E^{p,q}_2=H^p\bigl(X;K^q(\pt)\bigr)
\]
就是 Atiyah--Hirzebruch 谱序列的起点。

上述全部讨论不仅对 K 理论成立，对任何广义上同调理论也成立。在复 K 理论这一具体情形，事情还会进一步简化。得到
\[
E^{p,q}_2=\begin{cases}H^p(X;\Z) & q \text{ 偶},\\ 0 & q \text{ 奇},\end{cases}
\]
于是整个谱序列只需用 $q=0$ 那一行计算。$k$ 为偶数的微分 $d_k$ 全为零，因为它们交换奇偶 $q$。$k$ 为奇数时，微分为上同调运算：
\[
d_k:H^p(X;\Z)\to H^{p+k}(X;\Z).
\]
因此，计算复 K 理论从整系数上同调开始。然后依次对 $k$ 为奇数的运算 $d_k$ 取上同调。最后，把偶与奇上同调群剩余部分用一串（常常非平凡的）群扩张组合起来，还原 $K^0(X)$ 与 $K^1(X)$。这个论证的一个推论是：K 理论从不超过整系数上同调。

\transnote{「$q$ 偶时为 $H^p(X;\Z)$、$q$ 奇时为零」只在复 K 理论成立，依据是 $K^q(\pt)$ 随 $q$ 的奇偶取 $\Z$ 或 $0$。实 K 理论的 $KO^q(\pt)$ 以 $8$ 为周期，$E_2$ 页并非只有一行，$d_k$ 也不再只作用在一行上。}

\begin{example}
考虑实射影空间 $X=\R P^n$。它的整系数上同调为
\[
H^*(\R P^n;\Z)\iso\begin{cases}
\Z & \text{次数 }0,\\
\Z/2\Z & 0 \text{ 与 } n \text{ 之间的偶数次数},\\
\Z & n \text{ 为奇数时的次数 } n,\\
0 & \text{其余情形}.
\end{cases}
\]
次数 $0$ 的分量 $\Z$ 必须留存到 K 理论，因为它对应一点的上同调或 K 理论（这个 $\Z$ 确实不出现在约化上同调与约化 K 理论中），故所有微分 $d_k$ 在它上面为零。此外不存在非零同态 $\Z/2\Z\to\Z$，故最高次分量 $\Z$（$n$ 为奇数时）总落在微分 $d_k$ 的余核中。$\Z/2\Z$ 分量只出现在偶数次数，故所有奇数微分 $d_k$ 为零；偶数微分又总为零。于是 Atiyah--Hirzebruch 谱序列在 $E_2$ 页退化。由此 $K^*(\R P^n)$ 在次数 $0$ 有一个 $\Z$ 分量，$n$ 为奇数时在次数 $1$ 有一个 $\Z$ 分量，次数 $0$ 还有由一串 $\Z/2\Z$-扩张得到的挠分量。但这些信息没有完全确定 $\R P^n$ 的 K 理论：每个 $\Z/2\Z$-扩张都有两种可能形式：
\[
0\to\Z/2\Z\to\Z/4\Z\to\Z/2\Z\to0,
\]
或
\[
0\to\Z/2\Z\to\Z/2\Z\oplus\Z/2\Z\to\Z/2\Z\to0.
\]
事实上，$K(\R P^n)$ 用前者（见 [6, Proposition 2.7.7]）：
\[
K^0(\R P^n)\iso\Z\oplus\Z/2^k\Z,\qquad
K^1(\R P^n)\iso\begin{cases}\Z & n \text{ 奇},\\ 0 & n \text{ 偶},\end{cases}
\]
其中 $k=2^{\lfloor n/2\rfloor}$。一串 $\Z/2\Z$ 分次分量经由群扩张逐次组合，得到 $\Z/2\Z$、$\Z/4\Z$、$\Z/8\Z$ 等等。
\end{example}

\transnote{$d_k$ 的下标 $k$ 与 $K^0(\R P^n)\iso\Z\oplus\Z/2^k\Z$ 中的 $k$ 是两回事：前者是谱序列页码，后者是循环群的指数，此处 $k=2^{\lfloor n/2\rfloor}$。同一符号在相邻两处表示不同对象，是本节最易读错的地方。}

% ===== 原书 p.50 =====
\begin{thebibliography}{99}
\bibitem{1} J. F. Adams. \emph{Lectures on Lie groups}. W. A. Benjamin, Inc., New York-Amsterdam, 1969.
\bibitem{2} J. F. Adams. \emph{Stable homotopy and generalised homology}. Chicago Lectures in Mathematics. University of Chicago Press, Chicago, IL, 1995.
\bibitem{3} M. F. Atiyah. Algebraic topology and elliptic operators. \emph{Comm. Pure Appl. Math.}, 20:237--249, 1967.
\bibitem{4} M. F. Atiyah. Bott periodicity and the index of elliptic operators. \emph{Quart. J. Math. Oxford Ser. (2)}, 19:113--140, 1968.
\bibitem{5} M. F. Atiyah. Algebraic topology and operators in Hilbert space. In \emph{Lectures in Modern Analysis and Applications. I}, pages 101--121. Springer, Berlin, 1969.
\bibitem{6} M. F. Atiyah. \emph{K-theory}. Advanced Book Classics. Addison-Wesley Publishing Company Advanced Book Program, 1989.
\bibitem{7} M. F. Atiyah and R. Bott. On the periodicity theorem for complex vector bundles. \emph{Acta Math.}, 112:229--247, 1964.
\bibitem{8} M. F. Atiyah and F. Hirzebruch. Vector bundles and homogeneous spaces. In \emph{Proc. Sympos. Pure Math., Vol. III}, pages 7--38. American Mathematical Society, Providence, R.I., 1961.
\bibitem{9} M. F. Atiyah and I. M. Singer. The index of elliptic operators on compact manifolds. \emph{Bull. Amer. Math. Soc.}, 69:422--433, 1963.
\bibitem{10} M. F. Atiyah and I. M. Singer. The index of elliptic operators. I. \emph{Ann. of Math. (2)}, 87:484--530, 1968.
\bibitem{11} M. F. Atiyah and I. M. Singer. The index of elliptic operators. III. \emph{Ann. of Math. (2)}, 87:546--604, 1968.
\bibitem{12} R. Bott. The index theorem for homogeneous differential operators. In \emph{Differential and Combinatorial Topology (A Symposium in Honor of Marston Morse)}, pages 167--186. Princeton University Press, Princeton, N.J., 1965.
\bibitem{13} R. Bott and L. W. Tu. \emph{Differential forms in algebraic topology}, volume 82 of Graduate Texts in Mathematics. Springer-Verlag, New York, 1982.
\bibitem{14} H. Cartan and S. Eilenberg. \emph{Homological Algebra}. Princeton Landmarks in Mathematics. Princeton University Press, Princeton, N.J., 1999.
\bibitem{15} F. Hirzebruch. Lectures on K-theory. In J. F. Adams, editor, \emph{Algebraic Topology: A Student's Guide}, volume 4 of London Mathematical Society Lecture Notes Series, pages 223--238. Cambridge University Press, London, 1972. Mimeographed lecture notes of the AMS Summer Topology Institute, Seattle, 1963.
\bibitem{16} J. W. Milnor and J. D. Stasheff. \emph{Characteristic classes}. Princeton University Press, Princeton, N.J., 1974.
\bibitem{17} G. Segal. Equivariant K-theory. \emph{Inst. Hautes \'Etudes Sci. Publ. Math.}, 34:129--151, 1968.
\end{thebibliography}
